\documentclass[11pt]{article}
\usepackage[margin=24mm]{geometry}
\usepackage{amsmath,amssymb,amsthm,mathtools,mathrsfs,hyperref}
\newtheorem{theorem}{Theorem}
\newtheorem{lemma}[theorem]{Lemma}
\newcommand{\R}{\mathbb R}
\newcommand{\eps}{\varepsilon}
\allowdisplaybreaks
\title{Fixed testing functions at singular C1 contacts}
\author{}\date{23 October 2026}
\begin{document}\maketitle

A finite semidefinite construction needs testing functions that remain
fixed as the supporting Bell inequality varies. We construct such a space
for the reciprocal C1 contact family on an explicit source and testing
rectangle. The ratio $H=h/\sigma$ is unrestricted, and either event
weight $D$ or $\beta$ may vanish.

Two identities control the state couplings through the sum of complementary
source weights. They give a fixed positive diagonal Gram matrix with
uniformly small corrections. We then separate its source and testing
coefficients by one scalar identity. Including every retained square gives
220 testing functions for $\lambda>0$ and 157 for $\lambda\le0$.
Their union therefore supplies at most 377 functions for the complete
real-state parent on this rectangle. The other source domains and
projective testing charts remain necessary for a global pencil.

\section{The rectangle and source formulas}

Fix the rational constants
\[
 \eps=2^{-160},\qquad \delta=2^{-100},\qquad
 c_0=2^{-8},\qquad A_*=2^{30},\qquad
 \gamma={c_0\over7A_*}={1\over7\,2^{38}}.
\]
The source domain consists of
\[
 0<x,\sigma\le\eps,\quad 0<z\le1,\quad 0<b<1,
 \quad r=bz<1,\quad 0<h<1,
\]
with the original strict gaps and the two nonnegative event
coefficients specified below. Set
\[
 Q=1-b,\quad \theta=1-h,\quad H=h/\sigma,\quad e=xh/\sigma,
 \quad E=1-e>0,\quad k=\sqrt x,\quad q=\sqrt z,
\]
\[
 K=1+x,\quad R_s=1+xr,\quad M=R_s-xz\theta>0,
 \quad J=1+xh-\sigma(1-r)>0,
\]
\[
 P=\sigma+K+\sigma z(1+bx+2xh)
                  +\sigma bx^2hz^2+x^2h\theta z,
\]
\[
\begin{gathered}
 p={z(1+xh)E+(1-r)(1+xhz)\over2P},\qquad
 l={(1+\sigma r)\{eR_s(1+\sigma z)+ME\}\over2P},\\
 \kappa={K(1+\sigma r)(1+\sigma z)\over2MP}.
\end{gathered}
\]
The original coefficients are
\[
\begin{gathered}
 m=\kappa x^2zH=\kappa xze,\qquad
 C={Kz(1+\sigma r)E\over2(1-r)P},\\
 A=p/r,\quad \alpha=l/(zQ),\quad a=p/b,\quad \ell=l/Q,\\
 D={K(1+\sigma z)(1+\sigma r)-JP\over2JP}\ge0,\qquad
 \beta={K(1+\sigma r)(1+xhz)-\theta P\over2\theta P}\ge0.
\end{gathered}
\]
Here $H$ has no upper bound. Neither $b$, $Q$, $\theta$, $1-r$,
nor $z$ is bounded below. The testing rectangle is independently
\[
                  0\le u,w\le\delta.
\]
For signed angles $u_s,w_s$, set
$\alpha_s=\operatorname{sgn}(u_s)$ and
$\beta_s=\operatorname{sgn}(w_s)$, choosing the value $1$ at zero.
The exact testing gauge is
\[
 (u,w,X,Y,Z,W)=(|u_s|,|w_s|,\beta_sX_s,\alpha_sY_s,
                                \alpha_s\beta_sZ_s,W_s).
\]
It preserves the parent polynomial. Composing every testing function
below with this one fixed piecewise semialgebraic map extends the
same feature count to $|u_s|,|w_s|\le\delta$; no square root of a
negative angle is taken. State coordinates $X,Y,Z,W$ are unrestricted
real numbers; no state ratio is bounded.

Define
\[
 H_0=1+zQ,\quad B_s=P+xz(J-\sigma),
\]
\[
 \mathcal H=R_sH_0+z\{(1-r)xh(2+xr)
 -e[2+x(1+h+r\theta)]+Ke^2\}.
\]
Two exact source identities, obtained by multiplication, are
\begin{align}
 2\theta P\beta&=hB_s-\sigma H_0,\label{eq:BE3}\\
 2P(JD+x\theta\beta)&=\sigma^2\mathcal H.\label{eq:BE4}
\end{align}
The angular quantities used below are
\[
\begin{gathered}
 \nu=(2D-2\beta\theta+h)/\sigma,\quad
 \lambda=(8Dh+4\beta\theta^2+2h^2)/\sigma-4bp\sigma z>0,\\
 K_A=4\kappa H-\lambda,\quad d_G=4D/\sigma^2,
 \quad \rho=\lambda x/\sigma,\quad S=d_G+\rho,\\
 L_0=\lambda\sigma,\quad q_0=\lambda k,\quad
 c=\max(0,1-K_A),\quad a_0=k/\sigma.
\end{gathered}
\]
The letter $q_0$ here is a source quantity, not the constant
coefficient of a later scalar channel.

The complete reciprocal parent being certified is
\[
\begin{split}
 F_w={}&4m(W-wX+uY-uwZ)^2+4C(1+w^2)(uW-Y)^2\\
 &+4A(1+u^2)(1+w^2)Y^2
 +(4D+2)(uwW+uX-wY-Z)^2\\
 &+(4\alpha+2)(1+u^2)(Y-wZ)^2
 +(4\beta+2)(1+w^2)(X+uZ)^2\\
 &-2(1+u^2)(1+w^2)Z^2.
\end{split}
\]
All four state coordinates occur; they need not be normalized.

\section{Uniform source bounds, without an upper bound on $H$}

All inequalities in this section follow under the larger condition
$\eps\le2^{-22}$, so in particular on the chosen rectangle.
Directly from the positive formula for $P$ and the definitions,
\[
 1\le P\le1+5\eps,\quad |J-1|,|M-1|\le\eps,
 \quad 1\le B_s\le1+7\eps,\quad 1\le H_0\le2.
\]
Initially $0<e<1$ gives $0\le\mathcal H\le4$: nonnegativity
uses \eqref{eq:BE4}. Eliminate $\beta$ to obtain
\[
 eB_s=xH_0+\sigma\mathcal H-2PJD/\sigma.
\]
Consequently
\begin{equation}\label{eq:basic-bounds}
 e\le2x+4\sigma\le6\eps,\qquad
 |\mathcal H-H_0|\le32\eps,\qquad
 0\le D\le4\sigma^2,\quad 0\le\beta\theta\le1.
\end{equation}
For example, in $\mathcal H-H_0$ the terms are bounded by
$2\eps+3\eps+3e+2e^2$, which is at most $32\eps$ here.
The same positive formulas give
\begin{equation}\label{eq:p-bounds}
 |p-H_0/2|\le16\eps,\quad |l-1/2|\le8\eps,
 \quad 1/4\le p,l,\kappa\le2,
 \quad \kappa\le1,\quad m\le6\eps^2,
 \quad C(1-r)\le1.
\end{equation}
In particular $4a,4\ell\ge1$. The unbounded positive coefficients
$a,\ell,C$ will be retained in whole squares and will not be
estimated separately.

Exact rearrangement gives
\[
 \nu=2D/\sigma+[H_0-ez(J-\sigma)]/P,
\]
\[
 \lambda=2H+2e\theta z(J-\sigma)/P
                  -2\theta H_0/P+8Dh/\sigma-4bp\sigma z.
\]
Thus
\begin{equation}\label{eq:lambda-bounds}
 |\nu|\le3,\qquad |\lambda-2H|\le5,
 \qquad |L_0-2h|\le8\sigma.
\end{equation}

For completeness, the cancellation needed when $H$ is large is
polynomial. Write
\[
 N_\kappa=K(1+\sigma r)(1+\sigma z)-MP=xF_x+\sigma F_\sigma.
\]
Here $F_\sigma=r-1+\sigma rz$, and $F_x$ is the polynomial
$(N_\kappa-N_\kappa|_{x=0})/x$. Its coefficient $\ell^1$ norm
in $(x,\sigma,z,b,h,r)$ is $37$, as can also be checked by
expanding the displayed product. Since all six arguments lie in
$[0,1]$,
\[
 H(2\kappa-1)={eF_x+hF_\sigma\over MP},\qquad
 |H(2\kappa-1)-h(r-1)|\le462\eps.
\]
The estimate $|MP-1|\le8\eps$ suffices for the second inequality.
Combining this with \eqref{eq:lambda-bounds} more precisely yields
\begin{equation}\label{eq:KA-bounds}
 |K_A-\{2\theta H_0-2h(1-r)\}|\le2^{12}\eps,
 \qquad |K_A|\le5,\quad 0\le c\le6.
\end{equation}
If $c>0$, then $h\ge1/8$: for $h<1/8$ the expression in braces
is at least $3/2$, whereas the error is less than $1/4$.
The earlier elimination identity also gives
\[
 a_0^2(hB_s-\sigma H_0)\le\mathcal H.
\]
Therefore, whenever $c>0$,
\begin{equation}\label{eq:c-positive}
 a_0\le8,\qquad L_0\ge1/8,\qquad
 1/\lambda\le8\sigma.
\end{equation}
Every term below containing $a_0$ or $1/\lambda$ from the angular
subtraction is multiplied by $c$. On $c=0$ that subtraction is
identically zero; no uniform estimate on those unused factors is
required.

Put
\[
 R_0=-4\{2(C+D)(r-1)+2p+r-1\}.
\]
The positive expression for $C(1-r)$ and the formula for $p$ give
$|R_0|\le256(x+\sigma)\le512\eps$. Since
\[
 (x+\sigma)^2H=xe+2\sigma e+\sigma h\le2\eps,
\]
\eqref{eq:lambda-bounds} further gives
\begin{equation}\label{eq:R0-bounds}
 |R_0|\sqrt\lambda\le1024\sqrt\eps.
\end{equation}
Finally the exact identity
\[
 S={2\mathcal H\over P}+d_G(1-J+2hx)
       +{2eH_0\over P}-{2e^2z(J-\sigma)\over P}-4bpxz
\]
shows
\begin{equation}\label{eq:S-bounds}
 |S-2H_0|\le256\eps,\qquad
 1\le S\le5,\quad 0\le d_G,\rho\le5,
 \quad 0<L_0\le3,\quad 0<q_0\le4.
\end{equation}
Here $q_0^2=L_0\rho$; all bounds use $\lambda>0$ but no lower
bound on $\lambda$ or upper bound on $H$.

\section{A finite formula for the main Gram}

The state roots are
\[
\begin{gathered}
 N_x=X+uhZ,\quad N_y=(Y-bwZ)/q,\quad
 R=wZ/q-uqW,\quad \Gamma=\sigma Z-uwW,\\
 A_\theta=q(u^2-x)W,\quad T_\theta=kq(u-k)W,
 \quad J_C=N_y+bR,\quad J_\ell=N_y-QR,\quad
 \Gamma_*=\sqrt{d_G}\Gamma.
\end{gathered}
\]
Let
\[
 B_1=\sqrt{\lambda\sigma}\,uZ,\quad
 B_2=\sqrt\rho\,uwW,\quad
 B_3=\sqrt\lambda\,xqW,\quad
 B_4=\sqrt{\lambda x}\,wZ/q.
\]
When $WZ>0$, put
$G_*=\operatorname{sgn}(W)(B_1B_2B_3B_4)^{1/4}$ and
$d_i=B_i-G_*$. This is the positive-state convention followed,
when necessary, by changing the signs of all roots together.
When $WZ\le0$
put $d_i=B_i$. Equivalently, for every state,
\[
 G_*=\sqrt{\lambda x}\,\operatorname{sgn}(W)
                \sqrt{uw\max(WZ,0)},\qquad d_i=B_i-G_*.
\]
Define
\[
 U_*=\sqrt{q_0}\sqrt u\,q(u-k)W,
 \qquad K_*=\sqrt{q_0}\sqrt k\,q(u-k)W.
\]
Use the twelve-root column
\[
 g=(N_x,J_C,J_\ell,\Gamma_*,A_\theta,T_\theta,
                         d_1,d_2,d_3,d_4,U_*,K_*)^{\mathsf T}.
\]
The two radial changes are source-only and uniformly bounded:
\[
 R=J_C-J_\ell,\qquad N_y=QJ_C+bJ_\ell.
\]
Set
\[
 D_*={\rm diag}(2,1,1,1,1/2,1/2,1/2,1/2,1/2,1/2,1,1/2).
\]
We now define a polynomial quadratic form $\mathcal E(g)$, which
fully specifies the correction matrix; a cross term is divided by
two when placed in its two symmetric entries.

For its coefficients put $t=u^2$, $\xi=w^2$, $g_s=h-Q$ and
\[
\begin{gathered}
 R_x=-4\{2Db+(1+2m)g_s-2\beta\theta\},\quad B_x=8\kappa e,\\
 A_y=8(p-l)+4z\{(1+2m)g_s+2Dh+2p\sigma\},\\
 \chi=8(pb+lQ)-4h+4eh(2D+1)-8Dh\\
 \hspace{12mm}+4z\{4Dbh+4pb\sigma+(1+2m)g_s^2+2\beta\theta^2\},\\
 k_1=2(1+2m)g_s^2+8Dhb+8pb\sigma+4\beta\theta^2.
\end{gathered}
\]
Then
\[
\begin{split}
 E_{yy}&=2z\{(1+2m)u^2+(1+2D)w^2\},\\
 E_{RR}&=z\{8bD+8bp\sigma+4b\sigma(2C+2D+1)(r-1)
                                  +k_1u^2+2b^2(2D+1)w^2\},\\
 E_{yR}&=4z\{u^2[2Dh+2p\sigma+(1+2m)g_s]+b(2D+1)w^2\\
 &\hspace{25mm}+\sigma(2C+1)(r-1)+2D[1+\sigma(r-1)]+2p\sigma\}.
\end{split}
\]
The entire first normal is represented by the linear form
\[
 F_1={\sqrt{d_G}\over S}u\Gamma_*
 +{\sqrt\rho\over S}\{kd_1-ud_2+(u-k)d_3\}
                         -{\rho\over S}\sqrt\sigma T_\theta.
\]
The entire second normal is represented by
\[
\begin{split}
 F_2={}&{R_0\sqrt{d_G}\over S}qw\Gamma_*
 +{R_0\rho\sigma z\over S}R
 -{R_0\sqrt\rho\over S}(qw+z\sqrt\sigma)(d_2-d_3)\\
 &+\left\{{R_0q_0z\over S}
       +{R_0\sqrt\lambda\sqrt\rho\over S}qw\right\}T_\theta
                                          -8pquwF_1.
\end{split}
\]
These are exact identities $F_1=u\Gamma$ and
$F_2=qw(R_0-8pu^2)\Gamma$. In particular they keep both complete
state-normal couplings, not only their contact restrictions.

Write
\[
 \eps_C=2/\lambda+u+k,\qquad
 \eps_G=2a_0u+(u+k)/L_0.
\]
The correction is exactly
\begin{align*}
 \mathcal E={}&E_{yy}N_y^2+E_{RR}R^2+E_{yR}N_yR
                  -4\nu N_xF_1+J_CF_2\\
 &-4uqw(2m+2D+1)N_xN_y+uqwR_xN_xR
                 +qwB_xN_xA_\theta+uA_yN_yA_\theta+\chi uA_\theta R\\
 &-\tfrac12\rho\sigma R^2+\sqrt{q_0}\sqrt u\,U_*R
       -{cd_G\over S}\eps_C\sum_{i=1}^4d_i^2
       -{cd_G\over S}u^2R^2-{c\over S}\eps_G\Gamma_*^2\\
 &-{c\over S}\{\sqrt{a_0u}\,U_*+\sqrt{a_0k}\,K_*\}^2.
\end{align*}
Define $G$ by
$g^{\mathsf T}Gg=g^{\mathsf T}D_*g+\mathcal E$.
This is the actual main Gram after retaining the positive excess
of the $N_x,J_C,J_\ell$ squares, splitting $K_A+c\ge1$ equally
between $A_\theta,T_\theta$, and retaining their nonnegative
difference. The averaged circuit identity and angular identity
used in this step are precisely those of the printed all-$H$
proof at
\href{https://github.com/anubhav70747/q222-finite-sdp/blob/0c1ce2e15ed067a6845464e2468559cea5274f2c/latex/certificates/c1_short.tex#L7550-L7740}
{lines 7550--7740}. This definition therefore prints every main-Gram
entry through one quadratic polynomial instead of twelve rows of
expanded rational functions.

\section{Rational coefficient bounds and fixed radicals}

Put $v=\sqrt u$. Every entry of $G-D_*$ is a polynomial in
$(v,w)$, with source coefficients, supported on
\[
 \{1,w,w^2,v,v^2,v^2w,v^4,v^4w\}
 \subset\{v^\alpha w^\beta:0\le\alpha\le4,0\le\beta\le2\}.
\]
The following elementary bounds provide a reproducible coefficient
estimate:
\[
 |R_x|,|A_y|\le32,\quad |\chi|\le64,\quad
 |B_x|\le48\eps,\quad |k_1|\le16.
\]
The coefficients of $E_{yy}$ on $(u^2,w^2)$ are bounded by $(4,4)$;
those of $E_{RR}$ on $(1,u^2,w^2)$ by $(64\sigma,16,4)$;
those of $E_{yR}$ by $(32\sigma,12,8)$.
These follow term by term from \eqref{eq:basic-bounds}--\eqref{eq:p-bounds};
in the constant terms use $C(r-1)$, never $C$ separately.

Here is a single positive-coefficient majorant which verifies
all remaining estimates. In it $s=\sqrt\eps$, all twelve root
letters are formal nonnegative indeterminates, and
$n=r'=J_C+J_\ell$. Define
\[
 f_1=u\Gamma_*+sd_1+ud_2+(u+s)d_3+sT_\theta,
\]
\[
\begin{split}
 f_2={}&512s^2w\Gamma_*+512s^4r'
  +512s^2(w+s)(d_2+d_3)\\
 &+(2048s^2+1024sw)T_\theta+16uwf_1.
\end{split}
\]
The absolute coefficient polynomial of $\mathcal E$ is bounded by
\begin{align*}
 \mathcal Q={}&(4u^2+4w^2)n^2
 +(64s^2+16u^2+4w^2)r'^2
 +(32s^2+12u^2+8w^2)nr'\\
 &+12N_xf_1+J_Cf_2+16uwN_xn+32uwN_xr'
      +48s^2wN_xA_\theta+32unA_\theta+64uA_\theta r'\\
 &+\tfrac52s^2r'^2+2vU_*r'
      +6(16s^2+u+s)\sum_i d_i^2+6u^2r'^2
      +(144u+48s)\Gamma_*^2\\
 &+48uU_*^2+96vU_*K_*+48sK_*^2.
\end{align*}
Indeed \eqref{eq:S-bounds} bounds $\sqrt{d_G}/S$, $\sqrt\rho/S$,
$d_G/S$, $\rho/S$ by one, while \eqref{eq:c-positive} gives
$\eps_C\le16\sigma+u+k$ and $\eps_G\le24u+8k$ when needed.
The two small coefficients in $F_2$ use
$|R_0q_0|\le2048\eps$ and \eqref{eq:R0-bounds}.

Expansion, or simply evaluation at one of this positive-coefficient
polynomial, gives coefficient $\ell^1$ norm $8590$. Its constant
testing part has norm $5030$ and every term of that part is
divisible by $s$. Consequently every source coefficient of every
nonconstant testing monomial in $G$ has absolute value below
$2^{20}$, and the source matrix $G(0,0)$ satisfies
\begin{equation}\label{eq:matrixbounds}
 |G_{ij}(0,0)-(D_*)_{ij}|\le2^{20}\sqrt\eps=2^{-60}.
\end{equation}
These estimates hold before any ordering of the tensor product is
assumed. They bound pure-source coefficients in their own field.

For $r=12$ the following exact quadratic-form decomposition uses
only fixed directions:
\[
\begin{split}
 G={}&\sum_i\left(G_{ii}-{11\over48}\right)e_ie_i^{\mathsf T}\\
 &+\sum_{i<j}{1/48+G_{ij}\over2}(e_i+e_j)(e_i+e_j)^{\mathsf T}\\
 &+\sum_{i<j}{1/48-G_{ij}\over2}(e_i-e_j)(e_i-e_j)^{\mathsf T}.
\end{split}
\]
Write each of its $144$ scalar channels as
\[
 q(c,t)=q^{(0)}(c)+\sum_{j=1}^{7}a_j(c)b_j(t),
\]
where $b_j$ run through the seven nonconstant monomials
in the displayed support. By \eqref{eq:matrixbounds}, for every
channel
\[
 q^{(0)}\ge c_0=1/256,\qquad |a_j|\le A_*=2^{30}.
\]
In fact the off-diagonal channels have constant term at least
$1/96-2^{-61}$ and the diagonal ones at least $13/48-2^{-60}$.
On the independent testing rectangle,
$0\le b_j\le\sqrt\delta=2^{-50}<\gamma$.
Hence the exact separated identity
\[
 q=(q^{(0)}-c_0)+\frac12\sum_{j=1}^{7}
 \bigl[(A_*+a_j)(\gamma+b_j)+(A_*-a_j)(\gamma-b_j)\bigr]
\]
has nonnegative factors in their own labelled fields. It proves
the required tensor SOS without an appeal to entrywise positivity
in a tensor ring or to a point-dependent real congruence.

All twelve roots are source-linear combinations of the fixed
testing functions
\[
 \mathcal V=\{X,Y,Z,W,uZ,wZ,uW,uwW,u^2W,
                u^{3/2}W,u^{1/2}W,
                  \operatorname{sgn}(W)\sqrt{uw\max(WZ,0)}\}.
\]
The last function includes both state-sign cases and is zero
when $WZ\le0$. Changing all state signs changes all twelve roots'
signs and preserves the quadratic form.
Therefore a fixed main-Gram dictionary is
\[
 \mathcal V\ \cup\ \{f\sqrt{\gamma+b_j},f\sqrt{\gamma-b_j}:
                         f\in\mathcal V,\ 1\le j\le7\}.
\]
It has $180$ displayed functions. All radicals are testing-only;
the remaining coefficients are source-only square roots. The retained
positive squares are included explicitly below.

\section{The two zero faces are present, and $H$ is unrestricted}

No inverse of $D$, $\beta$, $b$, $Q$ or a state coordinate was used
in the coefficient estimates or Gram decomposition. At $D=0$
the root $\Gamma_*$ is zero and all formulas remain literal.
At $\beta=0$ the same formulas apply. Testing axes are included.
The denominators $S,M,P,J,\sigma,k$ used above are nonzero on
the stated source domain; their potentially dangerous combinations
were bounded explicitly rather than omitted on a boundary.

For explicit nonemptiness set $x=\sigma=t$, $z=1$, $b=1/2$,
$h=Ht$, with $0<t\le\eps$. The numerators $D_N=2JPD$ and
$\beta_N=2\theta P\beta$ are exactly
\[
\begin{split}
 D_N={}&(3-H)t^2+(3/4-6H)t^3+(H^2+H/2)t^4\\
       &+(-7H^2/2+H/4)t^5+(H^3-H^2/2)t^6,\\
 \beta_N={}&(H-3/2)t+4Ht^2-Ht^3+4H^2t^4
                                      +(-H^3+H^2/2)t^5.
\end{split}
\]
Uniformly for $1\le H\le4$, the remainders beyond the first
terms have absolute values at most $(687/4)t^3$ and $156t^2$.
Thus continuity gives a $\beta=0$ source with $H\in[1,2]$ and
$D>0$, and a $D=0$ source with $H\in[2,4]$ and $\beta>0$.
On the first face $p\le7/8$ and
$\lambda=8DH+2(H^2-p)t\ge t/4$.
On the second, $\beta\theta\ge t/16$, so
$\lambda\ge1/8-4t>0$. All strict native gaps hold.

Finally take $\sigma=t$, $x=t^2/16$, $z=1$, $b=1/2$, $h=3/4$.
Then $H=3/(4t)$ is unbounded along $t\downarrow0$ and
\[
 D_N={-3t^2\over262144}
 (6t^5-61t^4+608t^3-7872t^2+8192t-126976)>0,
\]
\[
 \beta_N={3(3t^2+64)(2t^3+5t^2+64)\over16384}>0.
\]
Also $K_A=O(\eps)$ on this path by \eqref{eq:KA-bounds}, so
$c>0$ is represented. These are singular limits inside the
evaluated rectangle, not merely a nonsingular compact slice.


% Include after the positive-angular main-Gram construction.
% Requires amsmath, amssymb, amsthm and mathrsfs; theorem is already defined.
\section{Fixed testing functions for the complete positive-angular parent}
Retain the preceding source/testing rectangle and $\lambda>0$.
Compose the testing functions with the angle gauge already defined, then
multiply all four state coordinates by $\operatorname{sgn}_+(W)$,
where $\operatorname{sgn}_+(0)=1$. Thus $W\ge0$.
Put $Z_+=\max(Z,0)$ and $Z_-=\max(-Z,0)$. The common state sign
preserves the quadratic parent. Every testing function below is pulled
back through this one fixed map.

\subsection{The exact whole-parent decomposition}

All symbols $g,G,N_x,N_y,R,\Gamma,A_\theta,T_\theta,J_C,J_\ell$
retain their preceding definitions.
Set $F_A=u(Z-a_0wW)$, $X_e=X+ewW$, and
\[
 a_C=cd_G/S,\quad q_C^{(0)}=\tfrac12-a_C(2/\lambda+k),\quad
 p_C=q_C^{(0)}-a_C\delta.
\]
The proved bounds give $0\le a_C\le6$ and $p_C>1/4$.
The circuit remainder satisfies
\[
 \mathscr C=\sum_{i=1}^4d_i^2+\mathscr C_+,
\]
and the angular difference is
\[
 \Delta_A=(2/\lambda+u+k)\mathscr C+u^2R^2
       +\{2a_0u+(u+k)/L_0\}\Gamma^2-A_\theta^2.
\]
The complete parent is exactly
\begin{align*}
 F_w={}&g^{\mathsf T}Gg+\mathcal E_+
 +(2+4m)w^2N_x^2
 +\{4a+2z(2C+1)-1\}J_C^2+(4\ell-1)J_\ell^2\\
 &+\tfrac12L_0F_A^2
 +{K_A+c-1\over2}(A_\theta^2+T_\theta^2)
 +{K_A+c\over2}(A_\theta^2-T_\theta^2)\\
 &+{cd_G\over S}\Delta_A
 +\{p_C+a_C(\delta-u)\}\mathscr C_+,
\end{align*}
where
\[
\begin{split}
 \mathcal E_+={}&4\beta(1+w^2)N_x^2+2(2D+1)(uX_e)^2
                 +4bpzu^2\Gamma^2\\
 &+4a(u^2+w^2+u^2w^2)J_C^2+4\ell u^2J_\ell^2
                                             +4Czw^2J_C^2.
\end{split}
\]
After the expansions below, every square coefficient outside $G$
is nonnegative and source-only, except the explicitly separated
testing factor $\delta-u$, which is nonnegative on the fixed
testing rectangle. This equality
is the full-state identity, not its restriction to a stationary
state or a determinant. No inverse of $D$ or $\beta$ occurs.

\subsection{Ordinary retained squares: seven additional functions}

Let $\mathcal V$ be the twelve pre-Gram functions
\[
 \{X,Y,Z,W,uZ,wZ,uW,uwW,u^2W,
                     u^{3/2}W,u^{1/2}W,\sqrt{uwWZ_+}\}.
\]
All ordinary retained roots are source-linear combinations of
$\mathcal V$ and the seven functions
\[
                   wX,\ uwZ,\ uX,\ u^2wW,\ uY,\ wY,\ uwY.
\]
Here are all seventeen roots; every displayed square root of a
coefficient is source-only:
\[
\begin{array}{c|l}
\text{family}&\text{square roots}\\\hline
\mathcal E_+&2\sqrt\beta N_x,\quad2\sqrt\beta\,wN_x,\quad
 \sqrt{2(2D+1)}\,uX_e,\quad2\sqrt{bpz}\,u\Gamma\\
 &2\sqrt a\,uJ_C,\quad2\sqrt a\,wJ_C,\quad2\sqrt a\,uwJ_C,\quad
 2\sqrt\ell\,uJ_\ell,\quad2\sqrt{Cz}\,wJ_C\\
N_x\text{ excess}&\sqrt{2+4m}\,wN_x\\
\text{radial excess}&\sqrt{4a+2z(2C+1)-1}\,J_C,\quad
                              \sqrt{4\ell-1}\,J_\ell\\
\text{averaged square}&\sqrt{L_0/2}\,F_A\\
\text{angular excess}&\sqrt{(K_A+c-1)/2}\,A_\theta,\quad
                       \sqrt{(K_A+c-1)/2}\,T_\theta\\
\text{angular difference}&\sqrt{(K_A+c)/2}\,qu(u-k)W,\\
 &\sqrt{(K_A+c)/2}\,q\sqrt{2k}\sqrt u(u-k)W.
\end{array}
\]
For example $J_C=Y/q-bquW$ and
$J_\ell=Y/q-wZ/q+QquW$, so multiplying by $u,w,uw$ adds
exactly the stated functions. All other rows follow by substitution.
The last two squares use
$A_\theta^2-T_\theta^2=z(u-k)^2W^2(u^2+2uk)$.

\subsection{The angular difference: nineteen functions}

The following calculation clears the intermediate state ratio
completely. Define positive source constants
\[
 A_u=1/k,\quad A_v=1/(\sqrt\sigma q),\quad
 A_z=\sqrt\sigma/(kq),\quad M_0=k^2q,\quad a_*=\lambda k,
\]
and write $r_* =\sqrt{wW|Z|}$. Eight cleared roots are
\[
\begin{split}
 f_1={}&M_0(A_uA_vA_zuw|Z|-W),\\
 f_2={}&\sqrt2 M_0(A_vA_zw|Z|-A_uuW),\\
 f_3={}&M_0\{(\sqrt{a_*}+1/\sqrt{a_*})A_uA_vuwW
                   -\sqrt{a_*}A_uA_zu|Z|-A_z|Z|/\sqrt{a_*}\},\\
 f_4={}&\sqrt{A_u}\sqrt u\,f_3,\\
 f_5={}&\sqrt{a_*}M_0(W-A_vA_zw|Z|),\qquad
 f_6=\sqrt{A_u}\sqrt u\,f_5,\\
 f_7={}&\sqrt{2a_*A_vA_z}M_0(A_uu-1)r_*,\qquad
 f_8=\sqrt{A_u}\sqrt u\,f_7.
\end{split}
\]
For $Z\ge0$, $\Delta_A=\sum_{i=1}^8f_i^2$. To verify this,
put $U=A_uu$, $V=A_vw$, $Z_*=A_z|Z|/W$, $M=M_0W$;
the roots are $M$ times the eight roots obtained by expanding
every $U+1$ factor in the identity
\begin{align*}
 &[2+a_*(U+1)]\Psi+U^2Q_*^2
       +[2U+(U+1)/a_*]R_*^2-T_*^2\\
 &=(T_*+UQ_*)^2+2Q_*^2
 +(U+1)(\sqrt{a_*}U(V-Z_*)-R_*/\sqrt{a_*})^2\\
 &\quad+a_*(U+1)(1-VZ_*)^2
                         +2a_*(U+1)VZ_*(U-1)^2,
\end{align*}
where $\Psi=1+U^2V^2+U^2Z_*^2+V^2Z_*^2-4UVZ_*$,
$Q_*=VZ_*-U$, $R_*=Z_*-UV$, and $T_*=U^2-1$.
It is a polynomial identity by multiplication.

For $Z<0$, add three squares. Put
\[
 C_2=4,\quad C_1=8a_*+8+4/a_*,\quad C_0=16+8a_*+4/a_*.
\]
Their roots are
\[
 M_0\sqrt{A_uA_vA_z}\,\boldsymbol1_{Z<0}
 \bigl(\sqrt{C_2}A_u u^{3/2}r_*,\quad
       \sqrt{C_1A_u}\,u r_*,\quad
       \sqrt{C_0}\sqrt u\,r_*\bigr).
\]
Indeed changing $Z_*$ to $-Z_*$ increases the left side by
$M^2UV|Z_*|(C_2U^2+C_1U+C_0)$. This also proves the exact
negative-state-sector coefficients. Every displayed source
factor is a real closed field square root in the source factor.

The resulting nineteen testing functions, including overlaps
with $\mathcal V$, are
\[
\begin{split}
 \mathcal A=\{&W,uw|Z|,w|Z|,uW,uwW,u|Z|,|Z|,
 u^{3/2}wW,u^{3/2}|Z|,\sqrt u|Z|,\sqrt uW,\sqrt u w|Z|,\\
 &ur_*,r_*,u^{3/2}r_*,\sqrt u r_*,
 \boldsymbol1_{Z<0}u^{3/2}r_*,
 \boldsymbol1_{Z<0}ur_*,\boldsymbol1_{Z<0}\sqrt u r_*\}.
\end{split}
\]
Multiply the eleven root formulas above by the pure-source factor
$\sqrt{cd_G/S}$ for their contribution to $F_w$.
Every cleared formula is defined at $W=0$ and $Z=0$. After
substituting squares for $u,w,W,|Z|$, the identities are polynomial;
therefore they include those literal boundaries, without a limit
or a state-coordinate inverse in the final formulas.

\subsection{The circuit remainder: fourteen functions}

Set the positive source constants
\[
 s_1=\sqrt{\lambda\sigma},\quad s_2=\sqrt\rho,\quad
 s_3=\sqrt\lambda\,xq,\quad s_4=\sqrt{\lambda x}/q,
 \qquad g_0=\sqrt{\lambda x}.
\]
Define seven testing functions
\[
\begin{array}{ll}
 t_1=u^{3/4}w^{1/4}W^{1/4}Z_+^{3/4},&
 t_2=u^{3/4}w^{3/4}W^{3/4}Z_+^{1/4},\\
 t_3=u^{1/4}w^{1/4}W^{3/4}Z_+^{1/4},&
 t_4=u^{1/4}w^{3/4}W^{1/4}Z_+^{3/4},\\
 t_5=u^{3/4}w^{1/2}W^{1/2}Z_+^{1/2},&
 t_6=u^{1/4}w^{1/2}W^{1/2}Z_+^{1/2},\\
 t_7=\sqrt{uwWZ_-}.&
\end{array}
\]
An explicit full-sector SOS is
\[
 \mathscr C_+=r_1^2+r_2^2+r_3^2+r_4^2,
\]
\[
\begin{gathered}
 r_1=\sqrt{2g_0}(\sqrt{s_1}t_1-\sqrt{s_2}t_2),\quad
 r_2=\sqrt{2g_0}(\sqrt{s_3}t_3-\sqrt{s_4}t_4),\\
 r_3=2\sqrt{g_0}\{(s_1s_2)^{1/4}t_5-(s_3s_4)^{1/4}t_6\},
 \qquad r_4=2\sqrt{\lambda x}\,t_7.
\end{gathered}
\]
The first three roots vanish for $Z\le0$ and the last vanishes
for $Z\ge0$. They also vanish wherever $WZ=0$. On $Z>0$ they
are precisely the separated two-pairing four-monomial identity;
on $Z<0$ the remainder is $4\lambda xuwWZ_-$.

There is no point-dependent scalar reserve at this step:
\[
 (p_C+a_C(\delta-u))\mathscr C_+
      =\sum_{j=1}^4(\sqrt{p_C}r_j)^2
        +\sum_{j=1}^4(\sqrt{a_C}\sqrt{\delta-u}\,r_j)^2.
\]
Thus the fixed circuit feature list is exactly the seven $t_i$
and their seven products with $\sqrt{\delta-u}$.

\subsection{The complete testing space}

For the main Gram take the seven nonconstant monomials
\[
 \mathcal B=\{w,w^2,\sqrt u,u,uw,u^2,u^2w\},\qquad
 \gamma_7={1\over7\,2^{38}}.
\]
The main list is $\mathcal V$ times the fifteen fixed factors
$1,\sqrt{\gamma_7+b},\sqrt{\gamma_7-b}$ for $b\in\mathcal B$,
giving $180$ functions. Its $144$ fixed directions are $e_i$,
$e_i+e_j$, and $e_i-e_j$, with the exact scalar-channel formulas
derived above, using $7A_*\gamma_7=c_0$.

\begin{theorem}
On the stated independent rectangle, the full real-state parent
$F_w$ is a sum of squares of source-linear combinations of one
fixed list of $220$ labelled testing functions: the $180$ main
functions, the seven ordinary functions, the nineteen angular
functions, and the fourteen circuit functions. Repeated functions
are harmless; the real dimension of their span is at most $220$.
There is an explicit construction using at most $2196$ square
roots: $144\cdot15$ main roots, seventeen ordinary roots, eleven
angular roots, and eight circuit roots. The same list, composed
with the fixed gauges above, includes all real state-sign sectors,
both testing axes, and the source faces $D=0$ and $\beta=0$.
\end{theorem}

This theorem follows from the displayed exact decomposition and
the separated coefficient formulas. It needs no Gram factorization
with mixed radicals. Without using the physical reality reduction,
apply the theorem separately to the real and imaginary state
amplitudes; their union has at most $440$ functions and $4392$
roots. The parent is real quadratic in the state, so those two
contributions sum to its complex-state value.


\section{The complete nonpositive-angular branch}\label{sec:fixed-c1-nonpositive}

The same source identities also control $\lambda\le0$. In this case,
a source-only shear replaces the angular circuit and gives a fixed
positive diagonal Gram matrix.

\subsection{Domain and source bounds}

Keep the same rational source formulas and strict gaps, and
now assume $\lambda\le0$. Thus
\[
 \begin{gathered}
 0<x,\sigma\le\eps=2^{-160},\quad 0<z\le1,\quad
 0<b<1,\quad r=bz<1,\quad 0<h<1,\quad D,\beta\ge0,
 \\ 0\le u,w\le\delta=2^{-100}.
 \end{gathered}
\]
State coordinates $W,X,Y,Z$ are unrestricted. Set $k=\sqrt x$,
$q=\sqrt z$, $Q=1-b$, $H=h/\sigma$, $H_0=1+zQ$,
$d_G=4D/\sigma^2$, $L=\lambda/\sigma$, and $K_c=1+L^2$.
The source identities BE3--BE4 give the following finite bounds:
\[
 H\le2,\quad h\le2\sigma,\quad -8\le L\le0,\quad
 |d_G-2H_0|\le200\eps,\quad
 |4\kappa H-2H_0|\le400\eps.                 \tag{N1}
\]
The already established common bounds used below are
\[
 1\le H_0\le2,\quad |\nu|\le3,\quad p\le2,\quad
 m\le6\eps^2,\quad D\le4\sigma^2,\quad
 |R_0|\le512\eps,\quad |B_x|\le48\eps,
\]
\[
 |R_x|,|A_y|\le32,\quad |\chi|\le64,\quad
 |k_1|\le16,\qquad 4a,4\ell\ge1.           \tag{N2}
\]
Here all coefficient symbols have their displayed definitions in the
preceding main-Gram construction, and $R_\gamma=R_0-8pu^2$.
The coefficient bounds for $E_{yy},E_{RR},E_{yR}$ on their respective
testing supports $(u^2,w^2)$ and $(1,u^2,w^2)$ are
\[
 (4,4),\qquad (64\sigma,16,4),\qquad
 (32\sigma,12,8).                         \tag{N3}
\]
These bounds retain $C(1-r)$ rather than estimating the possibly
unbounded positive coefficient $C$. Their source-formula proofs in the
shared calculation do not use $\lambda>0$; only the subsequent
positive-angular construction there uses that sign.

Here is a short derivation of (N1), using only the shared source
formulas and common bounds. Multiplying $\lambda\le0$ by $\sigma$
gives
\[
 8Dh+4\beta\theta^2+2h^2\le4bp\sigma^2z.
\]
Thus $H^2\le2bpz\le4$, $h\le2\sigma$, and
$\beta\theta^2/\sigma^2\le bpz\le2$. In particular
$\theta\ge1-2\eps$ and $\beta\theta/\sigma^2\le3$.
Also every positive term in $L$ can be discarded to give
$L\ge-4bpz\ge-8$.
For the remaining estimates it is enough to use $\eps\le1/2000$.
The positive formula for $P$, with $h\le2\sigma$, yields
$1\le P\le1+4\eps$, and $1\le B_s\le1+6\eps$.
BE3 is exactly
\[
 HB_s-H_0=2P\beta\theta/\sigma.
\]
Consequently
\[
 |H-H_0|\le
 \left\{\frac{4(1+4\eps)}{1-2\eps}+12\right\}\eps
 <18\eps.
\]
Use the common bound $|\mathcal H-H_0|\le32\eps$ and
$|PJ-1|\le6\eps$, $PJ\ge1-\eps$. BE4 gives
\[
 d_G=\frac{2\mathcal H}{PJ}
             -\frac{4x\beta\theta}{J\sigma^2}.
\]
The first term differs from $2H_0$ by at most
$88\eps/(1-\eps)$ and the second has absolute value at most
$12\eps/(1-\eps)$, so their sum is less than $200\eps$.
Finally $|M-1|\le\eps$ implies $|MP-1|\le6\eps$, while
$1\le K(1+\sigma r)(1+\sigma z)\le1+4\eps$.
The formula for $\kappa$ therefore gives
$|\kappa-1/2|\le5\eps/(1-\eps)<39\eps$.
Together with $H\le2$, this proves
\[
 |4\kappa H-2H_0|
 \le8(39\eps)+2(18\eps)<400\eps.
\]

\subsection{Exact complete quadratic correction}

Retain the uniform proof's roots
\[
\begin{gathered}
 N_x=X+uhZ,\quad N_y=(Y-bwZ)/q,\quad
 R=wZ/q-uqW,\quad \Gamma=\sigma Z-uwW,\\
 A_\theta=q(u^2-x)W,\quad V=uw(u^2-x)W,\quad
 J_C=N_y+bR,\quad J_\ell=N_y-QR,\quad X_e=X+ewW.
\end{gathered}
\]
In particular $R=J_C-J_\ell$ and $N_y=QJ_C+bJ_\ell$.
Introduce the source-only shear
\[
 T=\Gamma+(L/2)V,\qquad
 g=(N_x,J_C,J_\ell,T,V,A_\theta)^{\mathsf T},\qquad
 D_* =\operatorname{diag}(2,1,1,1,1,2).
\]
This shear is invertible with inverse $\Gamma=T-(L/2)V$;
both directions have coefficients bounded by four.
The source baseline has the exact decomposition
\[
\begin{split}
 2H_0\Gamma^2+L\Gamma V+(1+L^2)V^2
 &=T^2+V^2+(2H_0-1)\Gamma^2+\tfrac34L^2V^2,\\
 2H_0A_\theta^2&=2A_\theta^2+2(H_0-1)A_\theta^2.
\end{split}                                                   \tag{N4}
\]
Thus the principal baseline is the fixed diagonal $D_*$, with the
last terms in (N4) retained as nonnegative source squares. This is
a source-only change of roots, not a congruence depending on testing
parameters. Below all final directions are fixed in these six coordinates.

Define $\mathcal E$ by the following quadratic polynomial, substituting
$\Gamma=T-(L/2)V$, $R=J_C-J_\ell$, and $N_y=QJ_C+bJ_\ell$:
\[
\begin{split}
 \mathcal E={}&E_{yy}N_y^2+E_{RR}R^2+E_{yR}N_yR
 +(d_G-2H_0)\Gamma^2-4\nu uN_x\Gamma
 +qw(R_0-8pu^2)J_C\Gamma\\
 &-4uqw(2m+2D+1)N_xN_y+uqwR_xN_xR
 +qwB_xN_xA_\theta+uA_yN_yA_\theta\\
 &+(\chi+\lambda)uA_\theta R
 +(4\kappa H-2H_0)A_\theta^2+Lu^2\Gamma^2\\
 &-K_c\sigma u^2 A_\theta^2
 -K_c\sigma u(u^2-x)A_\theta R
 +K_c(u^2-x)V\Gamma .
\end{split}                                                    \tag{N5}
\]
This specifies every entry of $G$ by
$g^{\mathsf T}Gg=g^{\mathsf T}D_*g+\mathcal E$; cross coefficients
are divided by two between symmetric entries.

For completeness the entire retained part is
\[
\begin{split}
 \Pi={}&4\beta N_x^2+(4\beta+2+4m)(wN_x)^2
 +2(2D+1)(uX_e)^2+4bpz(u\Gamma)^2\\
 &+4a\{(uJ_C)^2+(wJ_C)^2+(uwJ_C)^2\}
 +4\ell(uJ_\ell)^2+4Cz(wJ_C)^2\\
 &+[4a+2z(2C+1)-1]J_C^2+(4\ell-1)J_\ell^2\\
 &+(2H_0-1)\Gamma^2+\tfrac34L^2V^2
 +2(H_0-1)A_\theta^2.
\end{split}                                                    \tag{N6}
\]
Every coefficient in (N6) is nonnegative in the source field.
No mixed radical occurs.

The exact parent identity is
\[
                     F_w=\Pi+g^{\mathsf T}Gg.             \tag{N7}
\]
Indeed the uniform proof gives
\[
 K_AA_\theta^2+\mathscr S
 =4\kappa H A_\theta^2+Lu^2\Gamma^2+L\Gamma V
                                      +\lambda uA_\theta R.
\]
Add $K_c$ times its exact zero identity
\[
 V^2-\sigma u^2A_\theta^2
 -u\sigma(u^2-x)A_\theta R+(u^2-x)V\Gamma=0.
\]
Splitting the radial and $N_x$ pivots and then using (N4) gives
precisely (N5)--(N7), including both complete Gamma normal terms.

\subsection{Evaluated coefficient bound and fixed directions}

All entries of $G-D_*$ are polynomials on the fixed testing support
\[
 \{1,w,w^2,u,uw,u^2,u^2w,u^3\}.                 \tag{N8}
\]
Here is a positive-coefficient majorant for the absolute coefficient
polynomial of (N5). Treat its root letters as formal nonnegative
variables, put $n=J_C+J_\ell$, and $C_*=T+4V$. Then
\[
\begin{split}
 \mathcal Q={}&(4u^2+4w^2)n^2
 +(64\eps+16u^2+4w^2)n^2
 +(32\eps+12u^2+8w^2)n^2\\
 &+200\eps C_*^2+12uN_xC_*
 +(512\eps w+16u^2w)J_CC_*\\
 &+16uwN_xn+32uwN_xn+48\eps wN_xA_\theta
 +32unA_\theta+65uA_\theta n\\
 &+400\eps A_\theta^2+8u^2C_*^2
 +65\eps u^2A_\theta^2
 +65\eps(u^3+u)A_\theta n
 +65(u^2+\eps)VC_* .
\end{split}                                                   \tag{N9}
\]
For the two replacement terms use $|\lambda|\le8\eps$,
$|\chi+\lambda|\le65$, $K_c\le65$, and $x,\sigma\le\eps$.
All other terms use (N1)--(N3). In particular
$|2m+2D+1|\le4$, a deliberately loose bound.

Evaluating this positive polynomial at one gives total coefficient
norm $10189<2^{14}$. Its constant-testing part has norm
$6109\eps$. The successive summands contribute
\[
32,336,208,5000,60,2640,32,64,48,64,130,400,200,65,260,650
\]
to the first total. The constant part consists of
$(256+128+5000+400+325)\eps$.
Consequently the deliberately larger rational bounds
\[
 |G_{ij}(0,0)-(D_*)_{ij}|\le2^{20}\eps,
 \quad\text{every nonconstant coefficient has magnitude }\le2^{20}
                                                               \tag{N10}
\]
hold in the source field.

For six coordinates use the exact fixed-direction decomposition
\[
\begin{split}
 G={}&\sum_i(G_{ii}-5/16)e_ie_i^{\mathsf T}\\
 &+\sum_{i<j}\frac{1/16+G_{ij}}2(e_i+e_j)(e_i+e_j)^{\mathsf T}
 +\sum_{i<j}\frac{1/16-G_{ij}}2(e_i-e_j)(e_i-e_j)^{\mathsf T}.
\end{split}                                                    \tag{N11}
\]
There are $36$ scalar channels. Let $b_1,\ldots,b_7$ be the
nonconstant monomials of (N8), fix
\[
 c_0=2^{-8},\qquad A_*=2^{30},\qquad
 \gamma=\frac1{7\,2^{38}},
\]
and write a channel as $q_0(c)+\sum_j a_j(c)b_j(t)$.
The diagonal constant channels are at least $11/16-2^{20}\eps$;
the off-diagonal ones are at least $1/32-2^{19}\eps$.
Thus $q_0\ge c_0$ and $|a_j|\le A_*$. Also
$0\le b_j\le\delta<\gamma$ throughout the independent testing box.
The exact separated identity
\[
 q_0+\sum_ja_jb_j=(q_0-c_0)
 +\frac12\sum_{j=1}^7
 \{(A_*+a_j)(\gamma+b_j)+(A_*-a_j)(\gamma-b_j)\}       \tag{N12}
\]
uses $7A_*\gamma=c_0$. Every factor on its right is nonnegative
in its own labelled field. Multiplying these squares by the fixed
linear directions in (N11) proves (N7) by an evaluated finite
tensor sum of squares, without an infinitesimal-residue argument
on a finite rectangle.

\subsection{A fixed complete testing space}

Every coordinate of $g$ is source-linearly spanned by
\[
 \mathcal V_0=\{X,Y,Z,W,uZ,wZ,uW,uwW,u^2W,u^3wW\}.       \tag{N13}
\]
For example $J_C=Y/q-buqW$, $J_\ell=Y/q-wZ/q+QuqW$,
and $T=\sigma Z-uwW+(L/2)(u^3wW-xuwW)$.
Thus the main Gram uses the $150$ displayed functions
\[
 \mathcal V_0\ \cup\
 \{f\sqrt{\gamma+b_j},\ f\sqrt{\gamma-b_j}:
                         f\in\mathcal V_0,\ 1\le j\le7\}.      \tag{N14}
\]
All retained roots in (N6) lie in the source-linear span of
$\mathcal V_0$ together with the seven additional functions
\[
       \{wX,uwZ,uX,u^2wW,uY,wY,uwY\}.                     \tag{N15}
\]
For instance $uX_e=uX+euwW$ and
$u\Gamma=\sigma uZ-u^2wW$. This proves a complete-parent
dictionary of $157$ displayed testing functions. Each source-only
linear combination obtained from (N11)--(N12) and (N6) has a
nonnegative source coefficient, absorbed by its own source square
root. The formulas give the finite coefficient prescription; no
successful global pencil or global feature bound follows.

For signed testing angles compose the entire dictionary with the
source-independent angle gauge defined above; this does
not change the count. The polynomial identities specialize literally
at $u=0$, $w=0$, and arbitrary zero state coordinates. Applying
the same real certificate to the real and imaginary parts treats
complex states (with the corresponding doubled real dictionary).

No inverse of $D$ or $\beta$ appears. The $\beta=0$ face is allowed
literally. In this branch (N1) actually implies
$d_G\ge2-200\eps>0$, so the $D=0$ face has empty intersection
with the stated box and $\lambda\le0$. Actual $D=0$ sources are
therefore supplied by the positive-$\lambda$ branch, not by this note.
This distinction is a consequence of the explicit bounds, not an
additional exclusion.

\subsubsection*{Nonempty sources, including the zero-beta face}
For a completely explicit source path set
\[
 x=\sigma=t,\qquad z=b=1-t^2,\qquad
 h=t-4t^2+Bt^3,\qquad 16\le B\le18,\quad0<t\le\eps.
\]
All strict native gaps hold, since $0<h\le2t$, $E=1-h>0$,
and $J,M\ge1-t>0$. Let $p_N=2Pp$, $D_N=2JPD$,
$\beta_N=2\theta P\beta$, and
\[
 N_L=4D_Nh+2\beta_N\theta J+2h^2JP
                              -2bp_N\sigma^2zJ=\sigma^2JP L.
\]
Direct polynomial multiplication, grouping by powers of $t$, gives
\[
\begin{split}
 |D_N-t^2|&\le27640t^3,\\
 |\beta_N-(B-17)t^3|&\le24459t^4,\\
 |N_L-(2B-38)t^3|&\le248856t^4.
\end{split}
\]
For these bounds take absolute coefficient sums in $(t,B)$ and
replace every $B^j$ by $18^j$ in the remainders. Thus $D>0$ and
$L<0$ throughout this path. At $B=18$, also $\beta>0$, giving
an entirely rational strict source for each rational $t$.
The signs of $\beta_N$ at $B=16,18$ are opposite. Polynomial
continuity therefore gives $B\in(16,18)$ with $\beta=0$, while
$D>0$ and $L<0$ remain strict. This verifies actual zero-beta
sources in the branch; it is not only a formal boundary allowance.

\end{document}
